\section{Further consequences}\label{sec:consequences}

The signed-prefix theorem gives bounds for other finite-dimensional norms
and, through matrix transference, for arrays whose rows can be permuted
separately.

\subsection{Finite-dimensional \texorpdfstring{$\ell_p$}{lp} bounds}

\begin{corollary}[Finite-dimensional $\ell_p$ bounds]\label{cor:lp-steinitz}
Let $1\le p\le\infty$, with $1/\infty=0$, and let $d,N\ge1$ be integers.
Every sequence $v_1,\ldots,v_N\in\R^d$ with $\lVert v_i\rVert_p\le1$
admits signs $\varepsilon_i\in\{-1,1\}$ such that
\[
 \max_{0\le k\le N}
 \left\lVert\sum_{i=1}^k\varepsilon_i v_i\right\rVert_p
 \le C d^{\max\{1/p,\,1-1/p\}}.
\]
If additionally $\sum_{i=1}^N v_i=0$, there is a permutation
$\pi\in\mathfrak S_N$ such that
\[
 \max_{0\le k\le N}
 \left\lVert\sum_{i=1}^k v_{\pi(i)}\right\rVert_p
 \le C d^{\max\{1/p,\,1-1/p\}}.
\]
Here $C$ is the same absolute constant as in
Theorems~\ref{thm:signing} and~\ref{thm:steinitz}, independent of $p,d,N$.
For $1\le p\le2$, both bounds have sharp dimension order $d^{1/p}$:
for every $d$ there is a unit-ball sequence for which every signing has
a prefix of norm at least $d^{1/p}$, and a finite indexed zero-sum
unit-ball family for which every permutation has a prefix of norm at
least $\tfrac13d^{1/p}$.
\end{corollary}

\begin{proof}
For $1\le p\le2$, the norm comparisons
\[
 \lVert x\rVert_2\le\lVert x\rVert_p
 \le d^{1/p-1/2}\lVert x\rVert_2
\]
allow direct application of Theorems~\ref{thm:signing}
and~\ref{thm:steinitz}, followed by the bound $C d^{1/p}$ on all prefixes.
For $2\le p\le\infty$, use
\[
 \lVert x\rVert_p\le\lVert x\rVert_2
 \le d^{1/2-1/p}\lVert x\rVert_p.
\]
Divide the input vectors by $d^{1/2-1/p}$, apply the corresponding
Euclidean theorem, and scale back to obtain $C d^{1-1/p}$.
Scaling preserves the zero sum, and in each case the one signing or
permutation controls all prefixes. At $p=\infty$ this gives $Cd$.

For the lower bounds, let $1\le p\le2$. The ordered standard basis
$e_1,\ldots,e_d$ has final signed sum of norm $d^{1/p}$ for every signing.
For the ordering problem, take $e_1,\ldots,e_d$ and $d$ copies of
$w=-\tfrac1d(1,\ldots,1)$. This indexed family sums to zero and lies in
the unit ball, since $\lVert w\rVert_p=d^{1/p-1}\le1$.
If $d\ge2$, consider any permutation and its prefix just after the
$k=\lfloor d/2\rfloor$-th basis vector. If $t$ copies of $w$ have
appeared, put $a=t/d\in[0,1]$. The prefix sum $s$ has $k$ coordinates
equal to $1-a$ and $d-k$ equal to $-a$. Since both $k$ and $d-k$ are
at least $d/3$, convexity gives
\[
 \lVert s\rVert_p^p
 =k(1-a)^p+(d-k)a^p
 \ge\frac d3\bigl((1-a)^p+a^p\bigr)
 \ge\frac d3\,2^{1-p}.
\]
Thus $\lVert s\rVert_p\ge\tfrac12(2d/3)^{1/p}
\ge\tfrac13d^{1/p}$, for every permutation. When $d=1$, the family
is $(1,-1)$ and either permutation has a prefix of norm one.
\end{proof}

\subsection{A colorful row-permutation consequence}

B\'ar\'any's matrix transference principle turns the signed-prefix bound
into a row-permutation bound for arrays. Each row may have its own
permutation, while the controlled sums aggregate the same number of
entries from every row.

\begin{corollary}[Euclidean colorful Steinitz bound]\label{cor:colorful-steinitz}
Let $d,k,n\ge1$ be integers and let $a_i^j\in\R^d$, for
$1\le j\le k$ and $1\le i\le n$, satisfy
\[
 \lVert a_i^j\rVert_2\le1,
 \qquad
 \sum_{j=1}^k\sum_{i=1}^n a_i^j=0.
\]
There is one tuple of permutations
$(\pi_1,\ldots,\pi_k)\in\mathfrak S_n^k$ such that
\[
 \max_{0\le m\le n}
 \left\lVert\sum_{j=1}^k\sum_{i=1}^m a_{\pi_j(i)}^j\right\rVert_2
 \le 2C\sqrt d,
\]
where $C$ is the absolute constant in Theorem~\ref{thm:signing}.
The same tuple works for every $m$ and may be chosen jointly from the
whole array. Only the global total sum is assumed to vanish; the bound
controls the displayed aggregate synchronous prefixes, not each row
separately.
\end{corollary}

\begin{proof}
For an arbitrary $k\times n$ array $B=(b_i^j)$ in the Euclidean unit
ball $B_2^d$, with no zero-sum condition, list its entries column by
column:
\[
 v_{(i-1)k+j}=b_i^j
 \qquad(1\le i\le n,\ 1\le j\le k).
\]
Apply Theorem~\ref{thm:signing} to this sequence of length $kn$. With
$\varepsilon_i^j=\varepsilon_{(i-1)k+j}$, every complete-column prefix
satisfies
\[
 \left\lVert\sum_{j=1}^k\sum_{i=1}^m\varepsilon_i^j b_i^j\right\rVert_2
 =
 \left\lVert\sum_{\ell=1}^{km}\varepsilon_\ell v_\ell\right\rVert_2
 \le C\sqrt d.
\]
This is the column-major reduction in the proof of
\cite[Lemma~2.1]{Barany2024Matrix}, with the present signed-prefix bound
in place of its dimension-linear bound. Thus B\'ar\'any's matrix signed
constant satisfies $V(B_2^d)\le C\sqrt d$, uniformly in $k,n$.
Here $V$ allows an entrywise signing of an arbitrary array, whereas the
row-permutation constant $U$ ranges over total-zero arrays and controls
the same aggregate prefixes. His transference theorem
\cite[Theorem~2.2]{Barany2024Matrix} gives $U(B_2^d)\le2V(B_2^d)$ for
every number of rows. Shorter auxiliary arrays may be padded with
trailing zero columns when using the bound uniformly in the number of
columns. Since $B_2^d$ is a centrally symmetric convex body, this gives
the asserted bound.
\end{proof}

Taking $k=1$ in the regular-simplex example in the introduction shows that no
$o(\sqrt d)$ bound can hold uniformly over these arrays. Thus the order
in the dimension is sharp; the numerical constant $2C$ is not asserted
to be optimal. The row permutations are existential, with no online or
efficient algorithm claimed.
